\chapter{Sécurité, tests, intégration continue et déploiement}
\label{chap:devsecops}

\begin{introchapitre}
Un outil de supervision a accès à tout l'hôte : il doit lui-même être sûr. Ce chapitre présente
les mesures de sécurité, la chaîne d'intégration continue, les tests, le déploiement sur un
serveur et l'évaluation de la solution.
\end{introchapitre}

\section{Sécurité de l'application}

\begin{table}[htbp]
\centering
\small
\begin{tabularx}{\linewidth}{L{3.3cm}L{5.2cm}Y}
\toprule
\textbf{Risque} & \textbf{Mesure} & \textbf{Vérification} \\
\midrule
Accès non autorisé à l'API & API fermée par défaut, jeton porteur, seules \code{POST /api/auth/login} et \code{GET /api/health} sont publiques & test automatique sur toutes les routes, figure~\ref{fig:controles} \\
Vol de mots de passe & PBKDF2-SHA512, 100\,000 itérations, sel aléatoire \cite{rfc8018} & test de connexion \\
Vol de jeton en base & seule l'empreinte SHA-256 du jeton est stockée & revue de code \\
Force brute & 5 échecs par 15 minutes et par couple (IP, utilisateur) & test automatique \\
Session oubliée & durée maximale 12 h, inactivité 2 h, révocation à la déconnexion & tests automatiques \\
Mot de passe initial & généré au premier lancement, changement obligatoire & test automatique \\
Secrets exposés & secrets en écriture seule : l'API ne renvoie jamais le webhook ni le mot de passe SMTP & test automatique \\
XSS, clickjacking & en-têtes CSP, \code{X-Frame-Options: DENY}, \code{nosniff} \cite{owasp} & test automatique \\
Injection via identifiant & identifiants de conteneur validés par expression régulière & test automatique \\
WebSocket ouvert & jeton exigé à la connexion & test automatique \\
\bottomrule
\end{tabularx}
\caption{Mesures de sécurité de l'application}
\label{tab:securite-app}
\end{table}

\section{Sécurité des conteneurs}

Le socket Docker donne un contrôle total de l'hôte. En version 0.1 il était monté
directement dans les conteneurs ; en version 0.2, aucun conteneur applicatif n'y a accès. Chacun
passe par un proxy \cite{socketproxy} qui n'autorise que les appels nécessaires
(tableau~\ref{tab:proxies}), sur un réseau interne sans port publié.

\begin{table}[htbp]
\centering
\small
\begin{tabularx}{\linewidth}{L{3.4cm}L{2.4cm}Y}
\toprule
\textbf{Proxy} & \textbf{Client} & \textbf{Appels autorisés} \\
\midrule
\code{docker-proxy-ro} & agent & lecture des conteneurs (liste, inspection, statistiques) \\
\code{docker-proxy-ops} & dashboard & journaux et redémarrage ; ni création, ni \code{exec}, ni suppression \\
\bottomrule
\end{tabularx}
\caption{Proxies d'accès à Docker}
\label{tab:proxies}
\end{table}

Les conteneurs applicatifs suivent les recommandations du NIST SP 800-190 \cite{nist800190}
et de Docker \cite{dockersecurity} :
\begin{itemize}
  \item exécution sous un utilisateur non privilégié (uid 1000) ;
  \item système de fichiers en lecture seule, seul \code{/tmp} est inscriptible en mémoire ;
  \item suppression de toutes les capacités Linux (\code{cap_drop: ALL}) et
        \code{no-new-privileges} ;
  \item images multi-étapes : les outils de compilation ne sont pas dans l'image finale ;
  \item dashboard publié uniquement sur \code{127.0.0.1} par défaut.
\end{itemize}

\begin{figure}[htbp]
\centering
\souscapture[0.99\linewidth]{cap-26-docker-ps.png}{Conteneurs en fonctionnement (\texttt{docker ps})}\\[0.6em]
\souscapture{cap-30-curl-401.png}{Appel sans jeton : 401}\label{fig:curl401}\hfill
\souscapture{cap-31-proxy-403.png}{Appel interdit par le proxy : 403}\label{fig:proxy403}
\caption{Conteneurs et vérification des contrôles d'accès}
\label{fig:controles}
\end{figure}

\section{Intégration continue}

Chaque envoi sur GitHub déclenche le pipeline de la figure~\ref{fig:ci}, exécuté par GitHub
Actions.

\diagramme[0.95\linewidth]{ci-pipeline.pdf}{Pipeline d'intégration continue}{fig:ci}

\begin{itemize}
  \item \textbf{agent} : installation des dépendances et exécution des tests pytest ;
  \item \textbf{server} : \code{npm ci}, tests et \code{npm audit} (échec si une vulnérabilité
        haute est trouvée) ;
  \item \textbf{ui} : vérification des types TypeScript et construction ;
  \item \textbf{images} : construction des deux images puis analyse par Trivy \cite{trivy} ;
        le pipeline échoue en cas de vulnérabilité critique ou haute corrigeable.
\end{itemize}

Deux mesures protègent la chaîne elle-même : les actions tierces sont épinglées par leur
empreinte SHA complète plutôt que par une étiquette modifiable \cite{ghhardening}, et
Dependabot \cite{dependabot} propose chaque semaine la mise à jour des dépendances.

Le pipeline a joué son rôle pendant le stage : Trivy a bloqué une image contenant des
vulnérabilités dans des paquets du système de base et \code{npm audit} a signalé une dépendance
vulnérable. Les deux ont été corrigés par mise à jour avant la fusion
(figure~\ref{fig:ci-captures}).

\begin{figure}[htbp]
\centering
\souscapture{cap-20-ci-vert.png}{Pipeline réussi}\hfill
\souscapture{cap-21-trivy-rouge.png}{Trivy bloquant une image vulnérable}
\caption{Intégration continue en pratique}
\label{fig:ci-captures}
\end{figure}

\section{Tests}

\subsection{Tests automatisés}

Le projet compte \textbf{51 tests automatisés} : 37 pour l'agent (pytest) et 14 pour le
serveur (\code{node:test}). Ils s'exécutent en quelques secondes, sur une base SQLite
temporaire, sans Docker.

\begin{table}[htbp]
\centering
\small
\begin{tabularx}{\linewidth}{L{4.4cm}cY}
\toprule
\textbf{Fichier} & \textbf{Tests} & \textbf{Ce qui est vérifié} \\
\midrule
\code{test_alert_manager.py} & 11 & attente avant déclenchement, pic court ignoré, pas de doublon, résolution automatique, alerte acquittée, règle supprimée, corrélation, échec de l'instantané sans blocage \\
\code{test_detectors.py} & 18 & opérateurs des seuils, validation de \code{rules.yaml}, anomalie après apprentissage, séries plates, saut durable, formules Docker, format Prometheus \\
\code{test_migrations.py} & 4 & base neuve, réexécution sans effet, mise à niveau d'une base v0.1, index unique \\
\code{test_retention_and_seed.py} & 4 & agrégation et purge, purge des alertes, initialisation unique des règles \\
\code{api.test.js} & 14 & routes fermées, connexion et limitation, sessions, secrets, règles, cycle de vie, outbox, WebSocket, en-têtes \\
\midrule
\textbf{Total} & \textbf{51} & \\
\bottomrule
\end{tabularx}
\caption{Répartition des tests automatisés}
\label{tab:tests}
\end{table}

\begin{figure}[htbp]
\centering
\souscapture{cap-24-pytest.png}{Agent : \texttt{pytest}}\hfill
\souscapture{cap-25-npm-test.png}{Serveur : \texttt{npm test}}
\caption{Exécution des tests}
\label{fig:tests}
\end{figure}

\subsection{Scénarios de validation de bout en bout}

\begin{table}[htbp]
\centering
\small
\begin{tabularx}{\linewidth}{L{4cm}L{4.6cm}Y}
\toprule
\textbf{Scénario} & \textbf{Action} & \textbf{Résultat attendu et obtenu} \\
\midrule
Charge CPU d'un conteneur & \code{docker run} d'un conteneur qui boucle (stress) & alerte d'anomalie en environ 15 s, alerte de seuil après 120 s, regroupées dans un même incident \\
Pic court & charge de 10 secondes & aucune alerte (abandon à l'état PENDING) \\
Fin de charge & arrêt du conteneur de test & résolution automatique et message de résolution \\
Webhook indisponible & URL Slack invalide & nouvel essai à 15, 30, 60 puis 120 s, puis échec enregistré \\
Acquittement & bouton \og Acquitter \fg{} & l'alerte reste ouverte, sans nouvelle notification, puis se résout seule \\
Accès sans session & \code{curl} sur \code{/api/alerts} & réponse 401 \\
Action Docker interdite & création de conteneur via le proxy & réponse 403 \\
\bottomrule
\end{tabularx}
\caption{Scénarios de validation manuelle}
\label{tab:scenarios}
\end{table}

\section{Déploiement sur un serveur}

SentinelOps a été installé sur une instance cloud Linux vierge, sans outil
d'infrastructure-as-code, afin de vérifier qu'une personne peut l'installer seule à partir du
guide (annexe~\ref{annexe:installation}). La figure~\ref{fig:deploiement} montre le déploiement
obtenu.

\diagramme[0.9\linewidth]{deploiement.pdf}{Diagramme de déploiement}{fig:deploiement}

Deux modes d'accès sont possibles :
\begin{itemize}
  \item \textbf{tunnel SSH} (recommandé) : le dashboard reste lié à \code{127.0.0.1} et aucun
        port supplémentaire n'est ouvert ; l'opérateur exécute
        \code{ssh -L 4000:localhost:4000 utilisateur@serveur} ;
  \item \textbf{port public} : \code{sen run --public}, avec une règle de pare-feu limitée à
        l'adresse IP de l'opérateur, de préférence derrière un proxy HTTPS.
\end{itemize}

\begin{figure}[htbp]
\centering
\souscapture[0.32\linewidth]{cap-27-instance-cloud.png}{Instance cloud}\hfill
\souscapture[0.32\linewidth]{cap-28-pare-feu.png}{Règles du pare-feu}\hfill
\souscapture[0.32\linewidth]{cap-29-dashboard-serveur.png}{Dashboard du serveur}
\caption{Déploiement sur un serveur cloud}
\label{fig:cloud}
\end{figure}

\section{Évaluation}
\label{sec:evaluation}

\subsection{Comparaison des deux estimateurs}

Pour justifier le choix de la médiane, les deux méthodes ont été appliquées à la même série :
30 mesures autour de 20~\% puis un saut durable à 75~\%, avec les mêmes paramètres
(fenêtre de 30, $k = 3$, plancher 2).

\begin{table}[htbp]
\centering
\small
\begin{tabular}{lcc}
\toprule
\textbf{Estimateur} & \textbf{Mesures signalées} & \textbf{Alerte (durée 15 s)} \\
\midrule
Moyenne et écart-type & 3 (15 s) & non déclenchée \\
Médiane et MAD & 15 (75 s) & déclenchée \\
\bottomrule
\end{tabular}
\caption{Comportement face à un saut durable}
\label{tab:estimateurs}
\end{table}

Avec la moyenne, l'anomalie fait elle-même grimper l'écart-type et cesse d'être signalée avant
la fin de la durée de 15 secondes : l'incident passe inaperçu. Avec la médiane, elle reste
signalée jusqu'à ce que la nouvelle valeur occupe la moitié de la fenêtre, puis devient la
nouvelle normale.

\subsection{Délais de détection et de notification}

Avec une cadence de 5 s et un passage du travailleur d'outbox toutes les 3 s, une anomalie passe
à FIRING en 15 à 20 s, une règle de durée $d$ en $d$ à $d+5$ s (60 à 65 s pour le CPU de l'hôte),
et la notification part au plus 3 s plus tard.

\subsection{Consommation de ressources}

\begin{table}[htbp]
\centering
\small
\begin{tabular}{lccc}
\toprule
\textbf{Conteneur} & \textbf{CPU moyen} & \textbf{Mémoire} & \textbf{Taille de l'image} \\
\midrule
agent & \afaire{x \%} & \afaire{x Mo} & \afaire{x Mo} \\
dashboard & \afaire{x \%} & \afaire{x Mo} & \afaire{x Mo} \\
proxies ($\times$2) & \afaire{x \%} & \afaire{x Mo} & \afaire{x Mo} \\
\bottomrule
\end{tabular}
\caption{Consommation de ressources (relevée avec \texttt{docker stats} et \texttt{docker images})}
\label{tab:ressources}
\end{table}

\subsection{Limites}

\begin{itemize}
  \item un seul hôte est supervisé par installation ;
  \item SQLite limite le volume de données et l'écriture à un processus à la fois ;
  \item le dashboard ne gère pas HTTPS lui-même : un proxy inverse est nécessaire en accès
        public ;
  \item un seul rôle (administrateur) ; pas de gestion fine des droits ;
  \item la livraison au moins une fois peut produire un doublon de notification après un échec
        partiel.
\end{itemize}

\begin{conclusionchapitre}
SentinelOps protège l'hôte qu'il supervise, vérifie chaque modification par des tests et une
analyse de sécurité automatiques, et s'installe sur un serveur réel en quelques minutes. Ses
limites, identifiées clairement, orientent les perspectives présentées en conclusion.
\end{conclusionchapitre}
